\chapter{Descripci\'o dels lliurables}
\label{chap:lliurables}

\section{Lliurable 1 --- Document LaTeX (20\%)}

L'alumne haur\`a de redactar un document formal d'entre \textbf{15 i 20 p\`agines} sobre un tema de la seva elecci\'o, formatat \'integrament en LaTeX. Es valorar\`a l'estructura global del document, la claredat expositiva, la qualitat del format, la coher\`encia interna i l'\'us adequat de les eines de LaTeX.

El document espec\'ific amb les instruccions detallades i la r\'ubrica d'avaluaci\'o es publicar\`a com a m\'inim \textbf{2 setmanes abans} de la data de lliurament.

\section{Lliurable 2 --- Torneig de Bots (25\%)}

L'alumne haur\`a de programar un \textbf{bot} en Python que competir\`a en un torneig estructurat en fases contra bots de refer\`encia i, posteriorment, contra els bots dels companys. El torneig constar\`a de \textbf{5 nivells de dificultat}. Per a cada nivell superat, l'alumne obtindr\`a \textbf{2 punts} sobre 10.

\textbf{Condici\'o de qualificaci\'o:} cal superar un m\'inim de \textbf{3 nivells} per obtenir qualsevol puntuaci\'o en aquest lliurable. No assolir aquest llindar implica una nota de \textbf{0} en aquest component.

Els alumnes que superin el llindar accediran a una \textbf{fase de torneig} contra els bots dels companys. En aquesta fase, \(K\) denota el conjunt d'alumnes classificats al torneig i \(W \subseteq K\) el \textbf{top 4}, que rep \textbf{+1}. Els alumnes que no pertanyen a \(K\) reben \textbf{0 punts extra}, i la resta de classificats segueixen una regressi\'o lineal estricta entre 0 i 1 en funci\'o de la seva posici\'o final. En l'exemple amb \(|K|=15\) i \(|W|=4\), la posici\'o 15 correspon a \(\tfrac{1}{12}\) i la posici\'o 5 a \(\tfrac{11}{12}\). Un cop en quedin quatre, es disputaran unes semifinals per fixar la classificaci\'o final.

El document espec\'ific amb les regles completes, el format del torneig i la r\'ubrica es publicar\`a com a m\'inim \textbf{2 setmanes abans} de la data de lliurament.

\section{Lliurable 3 --- Problemes Setmanals (25\%)}

Al llarg del curs, l'alumne haur\`a de resoldre i lliurar una s\`erie de \textbf{problemes setmanals} de programaci\'o. Aquests problemes es publicaran de manera peri\`odica i hauran de ser lliurats dins del termini indicat en cada cas.

El document espec\'ific amb el format, els criteris de lliurament i la r\'ubrica d'avaluaci\'o es publicar\`a com a m\'inim \textbf{2 setmanes abans} d'iniciar aquesta activitat.

\section{Lliurable 4 --- Projecte Final (30\%)}

El tema i el format del projecte final es \textbf{decidiran conjuntament amb els alumnes} a mesura que el curs avanci, en funci\'o del ritme assolit, dels interessos del grup i del grau d'autonomia adquirit.

Aquest lliurable podr\`a incorporar \textbf{+1 punt extra} addicional. Les condicions exactes d'obtenci\'o d'aquest bonus es definiran quan es publiqui el document espec\'ific del projecte.

Un cop definit el projecte, el document espec\'ific corresponent es publicar\`a com a m\'inim \textbf{2 setmanes abans} de la data de lliurament.
